\SourceLeaf{040}
\fancyhead[L]{\small Ein weiterer neuer Beweis}
\begin{center}
\textbf{EIN WEITERER NEUER BEWEIS DES SATZES,}\\[12pt]
\textbf{DASS JEDES POLYNOM}\\[12pt]
IN EINER VARIABLEN\\[12pt]
\textbf{\parbox{.95\linewidth}{\centering IN REELLE FAKTOREN ERSTEN ODER ZWEITEN GRADES\\ZERLEGT WERDEN KANN.}}\\[12pt]
\rule{3cm}{0.3pt}
\end{center}

\SecondArticle{1} Obwohl der Beweis des Satzes über die Zerlegung von Polynomen in Faktoren, den ich in einer vor sechzehn Jahren veröffentlichten Abhandlung gegeben habe, sowohl an Strenge als auch an Einfachheit nichts zu wünschen übrig zu lassen scheint, hoffe ich doch, dass es den Mathematikern nicht unwillkommen sein wird, wenn ich auf dieselbe höchst wichtige Frage zurückkomme und versuche, einen anderen, nicht weniger strengen Beweis aus ganz anderen Grundsätzen zu begründen. Jener frühere Beweis hängt nämlich wenigstens zum Teil von geometrischen Betrachtungen ab; der Beweis, den ich hier darzustellen beginne, wird dagegen auf rein analytischen Grundsätzen beruhen. In der angeführten Abhandlung habe ich die bedeutenderen analytischen Methoden besprochen, mit denen andere Mathematiker bis dahin unseren Satz zu beweisen versucht hatten, und ihre Mängel ausführlich dargelegt. Der schwerwiegendste und tatsächlich grundlegende dieser Mängel ist allen jenen Versuchen ebenso wie den neueren, die mir bekannt geworden sind, gemeinsam. Doch habe ich schon damals erklärt, dass er bei einem analytischen Beweis keineswegs unvermeidlich zu sein scheint. Die Sachkundigen mögen nun beurteilen, ob die damals gegebene Zusage durch diese neuen Bemühungen vollständig eingelöst ist.

\SecondArticle{2} Der Hauptuntersuchung werden einige Voruntersuchungen vorausgeschickt, sowohl damit nicht der Eindruck entsteht, etwas fehle, als auch weil ihre Behandlung vielleicht selbst auf Dinge, die andere bereits\SourceLeaf{041}
 berührt haben, noch etwas neues Licht werfen kann. Zunächst werden wir den größten gemeinsamen Teiler zweier Polynome in einer Unbestimmten behandeln. Dabei sei vorweg bemerkt, dass hier stets nur von Polynomen die Rede ist: Bildet man aus zweien von ihnen ein Produkt, so heißt jedes ein \textit{Teiler} dieses Produkts. Der \textit{Grad} eines Teilers bestimmt sich nach dem Exponenten der höchsten darin vorkommenden Potenz der Unbestimmten, ohne jede Rücksicht auf die Zahlenkoeffizienten. Im Übrigen können wir die Eigenschaften gemeinsamer Teiler von Polynomen umso kürzer behandeln, als sie denen der gemeinsamen Teiler von Zahlen vollkommen entsprechen.

Sind zwei Polynome $Y,Y'$ in der Unbestimmten $x$ gegeben, von denen das erste einen höheren oder wenigstens keinen niedrigeren Grad als das zweite hat, so bilden wir die folgenden Gleichungen:
\[
\begin{gathered} Y=qY'+Y''\\ Y'=q'Y''+Y'''\\ Y''=q''Y'''+Y''''\\ \text{usw. bis}\\ Y^{(\mu-1)}=q^{(\mu-1)}Y^{(\mu)}
\end{gathered}
\]
Dabei gilt folgende Regel: Zuerst wird $Y$ auf die übliche Weise durch $Y'$ dividiert; dann $Y'$ durch den Rest $Y''$ der ersten Division, der einen niedrigeren Grad als $Y'$ hat; dann wieder der erste Rest durch den zweiten, $Y'''$, und so fort, bis man zu einer Division ohne Rest gelangt. Dass dies schließlich notwendig eintreten muss, folgt daraus, dass der Grad der Polynome $Y',Y'',Y'''$ usw. fortwährend abnimmt. Dass diese Funktionen ebenso wie die Quotienten $q,q',q''$ usw. \textit{Polynome} in $x$ sind, braucht kaum erwähnt zu werden. Danach ist klar:

I. Geht man von der letzten dieser Gleichungen zur ersten zurück, so erkennt man, dass das Polynom $Y^{(\mu)}$ jedes der vorhergehenden teilt und daher gewiss ein gemeinsamer Teiler der vorgegebenen Polynome $Y,Y'$ ist.

II. Geht man von der ersten Gleichung zur letzten weiter, so zeigt sich, dass jeder gemeinsame Teiler der Polynome $Y,Y'$ auch jedes der folgenden und damit auch das letzte, $Y^{(\mu)}$, teilt. Die Polynome $Y,Y'$ können daher keinen gemeinsamen Teiler höheren Grades als $Y^{(\mu)}$ haben. Jeder gemeinsame Teiler desselben Grades wie $Y^{(\mu)}$ steht zu diesem im Verhältnis einer Zahl zu einer Zahl; daher ist dieser selbst als größter gemeinsamer Teiler anzusehen.
\par
\SourceLeaf{042}
III. Hat $Y^{(\mu)}$ den Grad 0, ist es also eine Zahl, so kann keine Funktion der Unbestimmten $x$ im eigentlichen Sinne die Polynome $Y,Y'$ teilen. In diesem Fall sagt man daher, diese Polynome hätten keinen gemeinsamen Teiler.

\Needspace{18\baselineskip}IV. Man nehme die vorletzte unserer Gleichungen und eliminiere daraus $Y^{(\mu-1)}$ mit Hilfe der drittletzten Gleichung; dann eliminiere man wiederum $Y^{(\mu-2)}$ mit Hilfe der vorhergehenden Gleichung und so fort. Auf diese Weise erhält man
\[
\begin{aligned}
Y^{(\mu)}&=+kY^{(\mu-2)}-k'Y^{(\mu-1)}\\
&=-k'Y^{(\mu-3)}+k''Y^{(\mu-2)}\\
&=+k''Y^{(\mu-4)}-k'''Y^{(\mu-3)}\\
&=-k'''Y^{(\mu-5)}+k''''Y^{(\mu-4)}\\
&\quad\text{usw.}
\end{aligned}
\]
wenn man annimmt, die Polynome $k,k',k''$ usw. seien nach der folgenden Regel gebildet:
\[
\begin{aligned}
k&=1\\
k'&=q^{(\mu-2)}\\
k''&=q^{(\mu-3)}k'+k\\
k'''&=q^{(\mu-4)}k''+k'\\
k''''&=q^{(\mu-5)}k'''+k''\\
&\quad\text{usw.}
\end{aligned}
\]
Damit gilt
\[
\pm k^{(\mu-2)}Y\mp k^{(\mu-1)}Y'=Y^{(\mu)}
\]
wobei die oberen Zeichen für gerades $\mu$ und die unteren für ungerades gelten. Haben also $Y$ und $Y'$ keinen gemeinsamen Teiler, so kann man auf diese Weise zwei Polynome $Z,Z'$ in der Unbestimmten $x$ finden, sodass
\[
ZY+Z'Y'=1
\]
V. Diese Aussage gilt offenbar auch umgekehrt: Kann die Gleichung
\[
ZY+Z'Y'=1
\]
so erfüllt werden, dass $Z,Z'$ Polynome in der Unbestimmten $x$ sind, dann können $Y$ und $Y'$ gewiss keinen gemeinsamen Teiler haben.

\SourceLeaf{043}
\Needspace{16\baselineskip}\SecondArticle{3}
Die zweite Voruntersuchung betrifft die Umformung symmetrischer Funktionen. Es seien $a,b,c$ usw. unbestimmte Größen, und ihre Anzahl sei $m$. Mit $\lambda'$ bezeichnen wir ihre Summe, mit $\lambda''$ die Summe ihrer Produkte aus je zweien, mit $\lambda'''$ die Summe ihrer Produkte aus je dreien usw., sodass die Entwicklung des Produkts
\[
(x-a)(x-b)(x-c)\, .\, .\, .\, .
\]
den Ausdruck
\[
x^m-\lambda'x^{m-1}+\lambda''x^{m-2}-\lambda'''x^{m-3}+\text{usw.}
\]
ergibt. Somit sind $\lambda',\lambda'',\lambda'''$ usw. symmetrische Funktionen der Unbestimmten $a,b,c$ usw., d.\,h. solche, in denen diese Unbestimmten auf dieselbe Weise vorkommen, oder deutlicher: solche, die bei jeder Vertauschung dieser Unbestimmten untereinander unverändert bleiben. Allgemeiner ist offenbar jedes Polynom in $\lambda',\lambda'',\lambda'''$ usw. (ob es nur diese Unbestimmten oder noch andere, von $a,b,c$ usw. unabhängige enthält) ein symmetrisches Polynom in den Unbestimmten $a,b,c$ usw.

\SecondArticle{4}
Der umgekehrte Satz liegt etwas weniger nahe. Es sei $\rho$ eine symmetrische Funktion der Unbestimmten $a,b,c$ usw.; sie besteht also aus einer bestimmten Anzahl von Gliedern der Form
\[
Ma^\alpha b^\beta c^\gamma\, .\, .\, .\, .
\]
wobei $\alpha,\beta,\gamma$ usw. nichtnegative ganze Zahlen bezeichnen und $M$ ein bestimmter oder wenigstens von $a,b,c$ usw. unabhängiger Koeffizient ist (falls außer $a,b,c$ usw. noch andere Unbestimmte in die Funktion $\rho$ eingehen). Zunächst legen wir für diese einzelnen Glieder eine bestimmte Ordnung fest. Dazu ordnen wir zuerst die Unbestimmten $a,b,c$ usw. in einer bestimmten, an sich völlig beliebigen Reihenfolge an, etwa so, dass $a$ die erste Stelle einnimmt, $b$ die zweite, $c$ die dritte usw. Von den beiden Gliedern
\[
Ma^\alpha b^\beta c^\gamma\, .\, .\, .\, .\quad\text{und}\quad Ma^{\alpha'}b^{\beta'}c^{\gamma'}\, .\, .\, .\, .
\]
weisen wir dann dem ersten eine höhere Stelle in der Ordnung zu als dem zweiten, wenn
\[
\begin{gathered}\text{oder}\quad\alpha>\alpha',\quad\text{oder}\quad\alpha=\alpha'\text{ und }\beta>\beta',\\\text{oder}\quad\alpha=\alpha',\ \beta=\beta'\text{ und }\gamma>\gamma',\quad\text{oder usw.}\end{gathered}
\]

\SourceLeaf{044}
also wenn die erste nicht verschwindende unter den Differenzen $\alpha-\alpha',\beta-\beta',\gamma-\gamma'$ usw. positiv ist. Da Glieder derselben Stelle in der Ordnung sich nur im Koeffizienten $M$ unterscheiden und daher zu einem Glied zusammengefasst werden können, nehmen wir an, dass die einzelnen Glieder der Funktion $\rho$ verschiedene Stellen in der Ordnung einnehmen.

Nun bemerken wir: Ist $Ma^\alpha b^\beta c^\gamma\, .\, .\, .\, .\, .$ unter allen Gliedern der Funktion $\rho$ das höchste in dieser Ordnung, so muss $\alpha$ größer oder wenigstens nicht kleiner als $\beta$ sein. Wäre nämlich $\beta>\alpha$, so hätte das Glied $Ma^\beta b^\alpha c^\gamma\, .\, .\, .\, .\, .$, das die Funktion $\rho$ wegen ihrer Symmetrie ebenfalls enthalten muss, eine höhere Stelle als $Ma^\alpha b^\beta c^\gamma\, .\, .\, .\, .\, .$, gegen die Voraussetzung. Ebenso ist $\beta$ größer oder wenigstens nicht kleiner als $\gamma$; ferner ist $\gamma$ nicht kleiner als der nächste Exponent $\delta$ usw. Daher sind die einzelnen Differenzen $\alpha-\beta$, $\beta-\gamma,\gamma-\delta$ usw. nichtnegative ganze Zahlen.

\Needspace{15\baselineskip}Zweitens betrachten wir ein Produkt aus beliebig vielen Polynomen in den Unbestimmten $a,b,c$ usw. Sein höchstes Glied muss das Produkt der höchsten Glieder dieser Faktoren sein. Ebenso ist klar, dass die höchsten Glieder der Funktionen $\lambda',\lambda'',\lambda'''$ usw. jeweils $a,ab,abc$ usw. sind. Daraus folgt, dass das höchste aus dem Produkt
\[
p=M{\lambda'}^{\alpha-\beta}{\lambda''}^{\beta-\gamma}{\lambda'''}^{\gamma-\delta}\, .\, .\, .\, .\, .
\]
hervorgehende Glied $Ma^\alpha b^\beta c^\gamma\, .\, .\, .\, .\, .$ ist. Setzt man also $\rho-p=\rho'$, so liegt das höchste Glied der Funktion $\rho'$ in der Ordnung gewiss niedriger als das höchste Glied der Funktion $\rho$. Offenbar sind aber $p$ und daher auch $\rho'$ symmetrische Polynome in $a,b,c$ usw. Behandelt man daher $\rho'$ ebenso wie vorher $\rho$, so zerfällt es in $p'+\rho''$, wobei $p'$ ein Produkt aus Potenzen von $\lambda',\lambda'',\lambda'''$ usw., multipliziert mit einem bestimmten oder wenigstens von $a,b,c$ usw. unabhängigen Koeffizienten, ist, während $\rho''$ ein symmetrisches Polynom in $a,b,c$ usw. ist, dessen höchstes Glied in der Ordnung niedriger liegt als das höchste Glied von $\rho'$. Fährt man auf dieselbe Weise fort, so wird $\rho$ offenbar schließlich auf die Form $p+p'+p''+p'''$ usw. gebracht, also in ein Polynom in $\lambda',\lambda'',\lambda'''$ usw. umgeformt.

\SecondArticle{5}
Den im vorigen Artikel bewiesenen Satz können wir auch so aussprechen: Ist irgendein symmetrisches Polynom $\rho$ in den Unbestimmten $a,b,c$ usw. gegeben, so lässt sich ein Polynom in ebenso vielen anderen Unbestimmten $l',l'',l'''$ usw. angeben, das durch die Substitutionen $l'=\lambda',l''=\lambda'',l'''=\lambda'''$ usw.
\SourceLeaf{045}
in $\rho$ übergeht. Außerdem zeigt man leicht, dass \textit{dies nur auf eine einzige Weise möglich ist.} Nehmen wir nämlich an, zwei verschiedene Funktionen der Unbestimmten $l',l'',l'''$ usw., etwa $r$ und $r'$, ergäben nach den Substitutionen $l'=\lambda',l''=\lambda'',l'''=\lambda'''$ usw. dieselbe Funktion von $a,b,c$ usw. Dann wäre $r-r'$ eine Funktion von $l',l'',l'''$ usw., die für sich nicht verschwindet, nach jenen Substitutionen aber identisch verschwindet. Dass dies widersinnig ist, erkennen wir leicht, wenn wir bedenken, dass $r-r'$ notwendig aus einer bestimmten Anzahl von Teilen der Form
\[
M{l'}^\alpha{l''}^\beta{l'''}^\gamma\, .\, .\, .\, .\, .
\]
besteht, deren Koeffizienten $M$ nicht verschwinden und die sich jeweils in ihren Exponenten unterscheiden. \Needspace{8\baselineskip}Die aus diesen einzelnen Teilen hervorgehenden höchsten Glieder werden daher durch
\[
Ma^{\alpha+\beta+\gamma+\text{usw.}}b^{\beta+\gamma+\text{usw.}}c^{\gamma+\text{usw.}}\, .\, .\, .\, .\, .
\]
dargestellt und gehören somit verschiedenen Stellen in der Ordnung an, sodass das höchste Glied überhaupt auf keine Weise wegfallen kann.

Die eigentliche Rechnung bei solchen Umformungen ließe sich durch mehrere Abkürzungen erheblich verkürzen. Dabei halten wir uns hier jedoch nicht auf, da für unseren Zweck schon die bloße Möglichkeit der Umformung genügt.
